\section{Antecedentes}
La implementación de maquetas en los laboratorios de Ingeniería de Control como lo es el sistema de control de nivel de 3 tanques de agua interconectados son herramientas de sumo valor para el aprendizaje de los estudiantes de ingeniería. Brindándoles la oportunidad de experimentar con un sistema físico en donde apliquen diversas disciplinas y metodologías de la teoría del control; como lo son el modelado matemático de sistemas físicos, el diseño de controladores que son vistos teóricamente y simulados en clase e inclusive la implementación de algoritmos de Identificación de Sistemas.\\\\

Pero para lograr todo esto, es necesario contar con un sistema de adquisición de datos que monitoree y almacene las variables del espacio de estado que describa al sistema.

\subsection{Sistemas de Adquisición Tradicionales}
    En la actualidad existen variadas formas de retribuir estas variables que los laboratorios de control utilizan para llevar a cabo sus experimentos:

     \begin{itemize}
        \item Tan sencillo como capturar a vista el registro de estos datos.
        \item Embebiendo sensores con frameworks amigables como Arduino.
        \item Herramientas que requieren de un conocimiento más avanzado: las tarjetas dedicadas a la adqusición de datos como los DAQs de National Instruments, NI ELVIS o el paquete de MATLAB Data Acquisition Toolbox. 
        \item Lenguajes de programación como LabVIEW y C/C++.
        \item Soluciones industriales como los PLCs.
     \end{itemize}

\subsection{Tecnologías Emergentes}
     Por otro lado, con tecnologías más recientes que ofrecen el desarrollo de software y hardware, han democratizado el desarrollo de sistemas embebidos, a tal grado que se pueden crear sistemas a la medida.

     \begin{itemize}
        \item \textbf{JavaScript:} Lenguaje de programación que permite la creación de aplicaciones web compatibles con la mayoría de los navegadores.
        \item \textbf{NodeJS:} Entorno de ejecución de JavaScript que permite la creación de aplicaciones del lado del servidor.
        \item \textbf{ReactJS:} Librería de JavaScript para la creación de interfaces gráficas de usuario.
        \item \textbf{SQLite:} Base de datos ligera que permite el almacenamiento de datos en un archivo local.
        \item \textbf{MQTT:} Protocolo de mensajería ligero que permite la comunicación entre dispositivos.
        \item \textbf{Modelo-Vista-Controlador (MVC):} Patrón de diseño de software.
        \item \textbf{Raspberry Pi:} Computadora de placa reducida que permite la creación de sistemas embebidos. \\ 
        Con enorme capacidad de procesamiento al grado de ejecutar múltiples tareas en paralelo  e interfáz de entradas y salidas digitales para propósito general.
     \end{itemize}

     Estas tecnologías, por su bajo costo, soporte para IoT y capacidad de alojar servidores web, las hacen ideales para integrar un sistema de adquisición de datos por red en conjunto al sistema de tanques compartidos.